% ================================================================
%  conteudo/objetivos.tex
% ================================================================

\section{Objetivos}

\subsection{Objetivo geral}

Avaliar o efeito da concentração inicial de substrato sobre a cinética de crescimento de \textit{Saccharomyces cerevisiae} em processo fermentativo conduzido em batelada, determinando os principais parâmetros cinéticos do processo.

\subsection{Objetivos específicos}

\begin{itemize}
    \item Monitorar o crescimento celular de \textit{S.\ cerevisiae} ao longo do tempo por meio da medição de absorbância (densidade óptica) e correlação com concentração de biomassa.
    \item Determinar os parâmetros cinéticos de crescimento: velocidade específica máxima de crescimento ($\mu_{max}$), tempo de geração ($t_g$) e fator de conversão substrato-biomassa ($Y_{X/S}$).
    \item Quantificar o consumo de glicose ao longo do processo fermentativo utilizando o método de DNS (ácido 3,5-dinitrossalicílico).
    \item Comparar os perfis cinéticos obtidos em diferentes concentrações iniciais de substrato e discutir os resultados à luz da literatura científica pertinente.
    \item Identificar as fases de crescimento microbiano (lag, exponencial, estacionária e declínio) nos perfis de biomassa obtidos experimentalmente.
\end{itemize}

